\section{Problem Statement}
\label{sec:problem}

\subsection{Preliminaries}

\fakeparagraph{Federated LoRA Fine-tuning}
We consider federated fine-tuning across $n$ clients, where client $i$ keeps its dataset $\mathcal{D}_i$ private and adapts a frozen LLM backbone $\phi$ through a LoRA module.
Let $\omega_i=|\mathcal{D}_i|/\sum_{k=1}^{n}|\mathcal{D}_k|$ denote the fraction of training examples held by client $i$, with $\sum_{i=1}^{n}\omega_i=1$.
In the standard homogeneous-rank setting, all clients optimize a shared adapter with factors $B\in\mathbb{R}^{d_{\mathrm{out}}\times r}$ and $A\in\mathbb{R}^{r\times d_{\mathrm{in}}}$.
With the LoRA scaling absorbed into the factors and the layer index omitted, the federated objective is
\begin{equation}
    \min_{B,A}
    \sum_{i=1}^{n}\omega_i
    \mathcal{L}_i\!\left(\phi+BA;\mathcal{D}_i\right),
    \label{equ:objective}
\end{equation}
where $\mathcal{L}_i$ is the mean training loss on client $i$'s local dataset.
Clients optimize local copies of the common factors, which the server coordinates toward the shared adapter $BA$.


\fakeparagraph{Federated LoRA Fine-tuning with Rank Heterogeneity}
In resource-heterogeneous deployments, client $i$ may instead operate under a fixed, client-specific rank budget $r_i$~\cite{bai2024flexlora,wang2024flora,cho2024hetlora}.
Under these client-specific feasible sets, the learning target becomes
\begin{equation}
    \begin{aligned}
        \min_{\{B_i,A_i\}_{i=1}^{n}}\quad
        &\sum_{i=1}^{n}\omega_i
        \mathcal{L}_i\!\left(\phi+B_iA_i;\mathcal{D}_i\right) \\
        \text{s.t.}\quad
        &B_i\in\mathbb{R}^{d_{\mathrm{out}}\times r_i},\quad
        A_i\in\mathbb{R}^{r_i\times d_{\mathrm{in}}},\quad i=1,\ldots,n,\\
        &\left|\{r_i:i=1,\ldots,n\}\right|>1.
    \end{aligned}
    \label{equ:heterogeneous-objective}
\end{equation}
The final constraint makes rank heterogeneity explicit by requiring heterogeneous client budgets.
Although the effective updates $B_iA_i$ share the common shape $d_{\mathrm{out}}\times d_{\mathrm{in}}$, their factor matrices have incompatible inner dimensions and therefore cannot be directly averaged across different ranks.
Consequently, federated knowledge transfer across heterogeneous adapters must preserve each client's rank budget while accounting for its local adaptation needs.

\fakeparagraph{Solution: Selective Rank Sharing}
To enable cross-client knowledge transfer without changing client $i$'s fixed budget $r_i$, we use a globally shared prefix of width $s$ and a client-specific protected remainder of width $p_i=r_i-s$ (see \figref{fig:selective-rank-sharing}).
Using subscripts $s$ and $p$ for the corresponding factor blocks, we write
\begin{equation}
    B_iA_i=B_{i,s}A_{i,s}+B_{i,p}A_{i,p},
    \qquad p_i=r_i-s,\quad s\in\{0,\ldots,\min_{1\le k\le n}r_k\}.
    \label{equ:selective-partition}
\end{equation}
Both blocks remain trainable.
After local training, the server aggregates the shared blocks separately for $B$ and $A$ using the same sample-size weights:
\begin{equation}
    \bar{B}_{s}=\sum_{i=1}^{n}\omega_i B_{i,s},
    \qquad
    \bar{A}_{s}=\sum_{i=1}^{n}\omega_i A_{i,s},
    \qquad
    \sum_{i=1}^{n}\omega_i=1.
    \label{equ:selective-aggregation}
\end{equation}
The server broadcasts $(\bar{B}_{s},\bar{A}_{s})$ to synchronize the shared prefix, while each client retains its protected block locally.
Selective rank sharing therefore controls how much of each heterogeneous adapter receives common external knowledge while preserving each client's total rank budget.





\subsection{Research Challenges}
\label{sec:problem:challenges}

Selective rank sharing preserves each client's fixed budget while allowing cross-client knowledge transfer, but it raises two coupled research challenges.
\begin{itemize}
    \item \textbf{How should a global receiving width be selected under heterogeneous rank budgets?}
    With a globally shared width $s$, client $i$ retains protected capacity $p_i=r_i-s$ within its fixed budget $r_i$.
    \figref{fig:acl-motivation}(a) shows that performance fluctuates as the shared rank changes and that the rank-16 and rank-64 groups respond differently on FLAN.
    At $s=8$, rank-16 clients lose ROUGE-L relative to $s=0$, while rank-64 clients gain.
    Moreover, the same width occupies different fractions of their budgets.
    These variations motivate dynamically selecting the global width $s$ rather than fixing it in advance.

    \item \textbf{How should the influence of external updates be adapted under coupled data and structural heterogeneity?}
    Clients differ in both their local data and their adapter rank.
    \figref{fig:acl-motivation}(b) shows that the relative similarity associated with matching rank versus matching domain differs across rank groups and training rounds.
    The rank-16 cosine gap remains positive from R1 onward, whereas the rank-64 gap changes sign across rounds.
    Thus, the relationship between client updates reflects both structural and data heterogeneity.
    Naive averaging may overlook client-dependent differences among contributions, motivating a \textit{fine-grained collaboration} mechanism.
\end{itemize}

Moreover, the two challenges are \textit{coupled}, as illustrated in \figref{fig:compound-heterogeneity}.
\begin{itemize}
    \item The allocation mechanism determines the subspace in which external influence may enter.
    This constrains the external knowledge available to the collaboration mechanism.
    \item The collaboration mechanism determines how strongly external knowledge is applied, which in turn affects the utility of the selected global width.
\end{itemize}
We therefore seek to coordinate both allocation and collaboration, so each client can adaptively absorb useful external knowledge while preserving its fixed budget and local adaptation needs.

\begin{figure}[t]
    \centering
    \begin{tikzpicture}[
        factor/.style={draw=blue!45!black, fill=blue!5, rounded corners=2pt,
            text width=4.1cm, minimum height=13mm, align=center, inner sep=5pt,
            font=\small},
        flow/.style={-{Stealth[length=2mm]}, thick, draw=gray!75}
    ]
        \node[factor] (allocation) at (0,0) {\textbf{Allocation}\\Select receiving and protected rank capacities};
        \node[factor] (collaboration) at (7.3,0) {\textbf{Collaboration}\\Adaptively weight external updates for each client};
        \draw[flow] ([yshift=2mm]allocation.east)
            to[bend left=18] node[midway, above, font=\scriptsize, fill=white, inner sep=1pt]
            {receiving capacity} ([yshift=2mm]collaboration.west);
        \draw[flow] ([yshift=-2mm]collaboration.west)
            to[bend left=18] node[midway, below, font=\scriptsize, fill=white, inner sep=1pt]
            {transfer utility} ([yshift=-2mm]allocation.east);
    \end{tikzpicture}
    \caption{Compound Impact of Allocation and Collaboration.}
    \label{fig:compound-heterogeneity}
\end{figure}
\FloatBarrier
